\documentclass[12pt]{article}
\usepackage[utf8]{inputenc}
\usepackage[T1]{fontenc}
\usepackage{hyperref}
\title{Rawas Language Resources for K-12 Curriculum Development}
\author{Compiled by Claude Sonnet 5.5 (claude-sonnet-5-5) for Andr\'as Kornai}
\date{}
\begin{document}
\maketitle

\section{Introduction}
The ISO 639-3 code rws (Rawas) is marked as retired. Glottolog lists it as a bookkeeping languoid (glottocode rawa1268). The code was merged into Musi [mui]. Rawas is treated as a variety of the Musi group of Malayic languages in South Sumatra, Indonesia, near the Musi Rawas area. The search found no stand-alone Rawas school materials. Any curriculum work would have to draw on Musi, Rejang and Malay materials.

\section{Primary Educational Resources}
No primers, textbooks or school programmes specific to Rawas were found in the search results.

\section{Linguistic and Reference Materials}
\begin{itemize}
\item Glottolog entry, ``Bookkeeping: Rawas (Retired)'': \url{https://glottolog.org/resource/languoid/id/rawa1268}
\item Wikipedia, Musi language: \url{https://en.wikipedia.org/wiki/Musi_language}
\item Wikipedia, Musi languages: \url{https://en.wikipedia.org/wiki/Musi_languages}
\item Robert McGinn (2005), ``What the Rawas Dialect Reveals about the Linguistic History of Rejang'', \emph{Oceanic Linguistics} 44(1). This is a historical study, not a general grammar. Listing: \url{https://signal.sciencespo-lyon.fr/article/512344/What-the-Rawas-Dialect-Reveals-about-the-Linguistic-History-of-Rejang}
\item Harun Mat Piah (1970), a Universiti Malaya thesis on a Rawa dialect in Perak, Malaysia. It describes a diaspora variety, not the Sumatran one. Glottolog reference: \url{https://glottolog.org/resource/reference/id/551504}
\end{itemize}

\section{Religious and Community Resources}
None found. The Asia Harvest people-group profile (\url{https://people-groups.asiaharvest.org/Indonesia/Rawas.pdf}) appeared in the results, but its figures are poorly sourced and it is not an educational resource.

\section{Digital and Audio Resources}
No audio, video or software resources specific to Rawas were found.

\section{Curriculum Development Context}
Rawas has no separate ISO status and no dedicated teaching materials in the search results. A curriculum would have to rest on new documentation work, using the Musi and Rejang literature as a starting point. The history of the Rawas river region is covered in \url{https://journal.ipm2kpe.or.id/index.php/KAGANGA/article/view/3826} (Dinamika Kehidupan Masyarakat di Sepanjang Sungai Rawas, 1864--1942). It is background reading only.

\end{document}